% Final acceptance tests (BGU \S 10.7.2.7).
\section{Final Acceptance Tests}
\label{sec:tests}

This section reports the measured performance of the delivered product
against every relevant metric. Every value comes from a frozen training run
recorded under \texttt{runs/final/} and aggregated by
\texttt{artifacts/Final\_Exp.json} (schema version 3). No values are
synthesized or interpolated.

\subsection{Test rig}
The campaign ran on a single workstation with the configuration in
Table \ref{tab:rig}. Determinism is guarded by
\texttt{pl.seed\_everything(42, workers=True)} for the canonical seed and by
the per-sample seed offset for the validation split.

\begin{table}[htbp]
\centering
\caption{Test rig and software stack.}
\label{tab:rig}
\setlength{\tabcolsep}{4pt}
\renewcommand{\arraystretch}{1.1}
\begin{tabular}{ll}
\toprule
Component & Value \\
\midrule
GPU                & NVIDIA RTX 5070 (Blackwell, sm\_120, 12 GB) \\
Precision          & \texttt{bf16-mixed} \\
CPU dataloader workers & 0 (Windows) \\
Framework          & PyTorch 2.11 \cite{paszke2019pytorch} + Lightning 2.x \cite{falcon2019lightning} \\
Image quality      & \texttt{torchmetrics} PSNR / SSIM, \texttt{data\_range = 1} \\
Determinism gate   & \texttt{src/tests/test\_degradation\_determinism.py} \\
\bottomrule
\end{tabular}
\end{table}

\subsection{Per-pipeline image quality}
\label{sec:tests:iq}

Table \ref{tab:pipelines} reports the mean PSNR and SSIM for every
degradation pipeline that appears in the experiments. The values are
identical across same-pipeline cells by construction, so Phase D rows
inherit the Phase B PSNR / SSIM at the same level. Figures
\ref{fig:samples_B_cifar} and \ref{fig:samples_C_cifar} show the visual
appearance of the pipelines on CIFAR-10; Figures \ref{fig:samples_B_mnist}
and \ref{fig:samples_C_mnist} repeat the comparison on MNIST.

\begin{table}[htbp]
\centering
\caption{Mean PSNR and SSIM per degradation pipeline. Each value is computed over a fixed deterministic 256-sample validation subset using torchmetrics with data\_range=1, sample seed offset = $2 \times 10^{7}$, clean reference upsampled to $224 \times 224$ by bilinear interpolation.}
\label{tab:pipelines}
\setlength{\tabcolsep}{4pt}
\renewcommand{\arraystretch}{1.1}
\begin{tabular}{llllrr}
\toprule
Dataset & Phase & Level & Axis & PSNR (dB) & SSIM \\
\midrule
CIFAR-10 & B & L1 & -- & 21.23 $\pm$ 1.54 & 0.801 $\pm$ 0.041 \\
CIFAR-10 & B & L2 & -- & 17.54 $\pm$ 1.32 & 0.697 $\pm$ 0.059 \\
CIFAR-10 & B & L3 & -- & 14.87 $\pm$ 1.41 & 0.635 $\pm$ 0.078 \\
CIFAR-10 & B & L4 & -- & 13.10 $\pm$ 1.41 & 0.601 $\pm$ 0.087 \\
CIFAR-10 & B & L5 & -- & 10.85 $\pm$ 1.81 & 0.566 $\pm$ 0.110 \\
CIFAR-10 & B2 & L1 & -- & 19.50 $\pm$ 2.31 & 0.816 $\pm$ 0.054 \\
CIFAR-10 & B2 & L2 & -- & 17.25 $\pm$ 1.88 & 0.733 $\pm$ 0.067 \\
CIFAR-10 & B2 & L3 & -- & 15.24 $\pm$ 1.72 & 0.673 $\pm$ 0.079 \\
CIFAR-10 & B2 & L4 & -- & 14.18 $\pm$ 1.86 & 0.652 $\pm$ 0.086 \\
CIFAR-10 & B2 & L5 & -- & 12.40 $\pm$ 2.50 & 0.624 $\pm$ 0.108 \\
CIFAR-10 & B2nr & L3 & -- & 15.24 $\pm$ 1.72 & 0.673 $\pm$ 0.079 \\
CIFAR-10 & C & L1 & Blur & 41.98 $\pm$ 2.54 & 0.983 $\pm$ 0.007 \\
CIFAR-10 & C & L1 & Noise & 28.60 $\pm$ 0.42 & 0.841 $\pm$ 0.021 \\
CIFAR-10 & C & L1 & Resolution & 30.78 $\pm$ 2.70 & 0.917 $\pm$ 0.035 \\
CIFAR-10 & C & L1 & S\&P & 21.66 $\pm$ 1.10 & 0.858 $\pm$ 0.023 \\
CIFAR-10 & C & L1 & Saturation & 36.45 $\pm$ 2.39 & 0.975 $\pm$ 0.010 \\
CIFAR-10 & C & L2 & Blur & 33.58 $\pm$ 2.59 & 0.952 $\pm$ 0.020 \\
CIFAR-10 & C & L2 & Noise & 23.34 $\pm$ 0.36 & 0.642 $\pm$ 0.038 \\
CIFAR-10 & C & L2 & Resolution & 25.33 $\pm$ 2.65 & 0.845 $\pm$ 0.060 \\
CIFAR-10 & C & L2 & S\&P & 18.61 $\pm$ 0.87 & 0.755 $\pm$ 0.032 \\
CIFAR-10 & C & L2 & Saturation & 31.34 $\pm$ 3.42 & 0.968 $\pm$ 0.014 \\
CIFAR-10 & C & L3 & Blur & 27.22 $\pm$ 2.54 & 0.885 $\pm$ 0.045 \\
CIFAR-10 & C & L3 & Noise & 20.12 $\pm$ 0.42 & 0.489 $\pm$ 0.047 \\
CIFAR-10 & C & L3 & Resolution & 21.75 $\pm$ 2.64 & 0.784 $\pm$ 0.079 \\
CIFAR-10 & C & L3 & S\&P & 16.38 $\pm$ 0.70 & 0.635 $\pm$ 0.035 \\
CIFAR-10 & C & L3 & Saturation & 27.94 $\pm$ 4.23 & 0.958 $\pm$ 0.022 \\
CIFAR-10 & C & L4 & Blur & 23.63 $\pm$ 2.52 & 0.822 $\pm$ 0.066 \\
CIFAR-10 & C & L4 & Noise & 17.87 $\pm$ 0.44 & 0.381 $\pm$ 0.050 \\
CIFAR-10 & C & L4 & Resolution & 19.42 $\pm$ 2.67 & 0.745 $\pm$ 0.091 \\
CIFAR-10 & C & L4 & S\&P & 14.96 $\pm$ 0.60 & 0.538 $\pm$ 0.038 \\
CIFAR-10 & C & L4 & Saturation & 25.41 $\pm$ 4.72 & 0.945 $\pm$ 0.033 \\
CIFAR-10 & C & L5 & Blur & 21.06 $\pm$ 2.52 & 0.772 $\pm$ 0.082 \\
CIFAR-10 & C & L5 & Noise & 15.49 $\pm$ 0.43 & 0.277 $\pm$ 0.048 \\
CIFAR-10 & C & L5 & Resolution & 15.88 $\pm$ 2.82 & 0.692 $\pm$ 0.109 \\
CIFAR-10 & C & L5 & S\&P & 13.88 $\pm$ 0.58 & 0.456 $\pm$ 0.041 \\
CIFAR-10 & C & L5 & Saturation & 24.17 $\pm$ 4.93 & 0.936 $\pm$ 0.041 \\
CIFAR-10 & C2 & L1 & Blur & 24.67 $\pm$ 5.66 & 0.943 $\pm$ 0.040 \\
CIFAR-10 & C2 & L1 & Resolution & 23.11 $\pm$ 3.97 & 0.883 $\pm$ 0.050 \\
CIFAR-10 & C2 & L1 & S\&P & 19.27 $\pm$ 2.05 & 0.824 $\pm$ 0.043 \\
CIFAR-10 & C2 & L2 & Blur & 23.74 $\pm$ 4.42 & 0.915 $\pm$ 0.044 \\
CIFAR-10 & C2 & L2 & Resolution & 21.22 $\pm$ 3.09 & 0.819 $\pm$ 0.068 \\
CIFAR-10 & C2 & L2 & S\&P & 17.19 $\pm$ 1.43 & 0.724 $\pm$ 0.044 \\
CIFAR-10 & C2 & L3 & Blur & 22.05 $\pm$ 3.35 & 0.855 $\pm$ 0.057 \\
CIFAR-10 & C2 & L3 & Resolution & 19.34 $\pm$ 2.65 & 0.763 $\pm$ 0.082 \\
CIFAR-10 & C2 & L3 & S\&P & 15.46 $\pm$ 1.03 & 0.609 $\pm$ 0.044 \\
CIFAR-10 & C2 & L4 & Blur & 20.43 $\pm$ 2.77 & 0.798 $\pm$ 0.071 \\
CIFAR-10 & C2 & L4 & Resolution & 17.79 $\pm$ 2.47 & 0.729 $\pm$ 0.092 \\
CIFAR-10 & C2 & L4 & S\&P & 14.28 $\pm$ 0.84 & 0.516 $\pm$ 0.044 \\
CIFAR-10 & C2 & L5 & Blur & 18.92 $\pm$ 2.46 & 0.753 $\pm$ 0.084 \\
CIFAR-10 & C2 & L5 & Resolution & 15.08 $\pm$ 2.49 & 0.681 $\pm$ 0.108 \\
CIFAR-10 & C2 & L5 & S\&P & 13.35 $\pm$ 0.74 & 0.437 $\pm$ 0.045 \\
MNIST & B & L1 & -- & 19.56 $\pm$ 1.75 & 0.577 $\pm$ 0.030 \\
MNIST & B & L2 & -- & 15.50 $\pm$ 1.30 & 0.423 $\pm$ 0.039 \\
MNIST & B & L3 & -- & 12.58 $\pm$ 1.30 & 0.318 $\pm$ 0.043 \\
MNIST & B & L4 & -- & 10.90 $\pm$ 1.31 & 0.248 $\pm$ 0.048 \\
MNIST & B & L5 & -- & 9.02 $\pm$ 1.96 & 0.247 $\pm$ 0.072 \\
MNIST & B2 & L1 & -- & 19.81 $\pm$ 1.56 & 0.864 $\pm$ 0.027 \\
MNIST & B2 & L2 & -- & 16.09 $\pm$ 1.45 & 0.737 $\pm$ 0.052 \\
MNIST & B2 & L3 & -- & 13.05 $\pm$ 1.54 & 0.586 $\pm$ 0.077 \\
MNIST & B2 & L4 & -- & 11.41 $\pm$ 1.38 & 0.490 $\pm$ 0.085 \\
MNIST & B2 & L5 & -- & 9.60 $\pm$ 2.42 & 0.419 $\pm$ 0.121 \\
MNIST & B2nr & L3 & -- & 13.05 $\pm$ 1.54 & 0.586 $\pm$ 0.077 \\
MNIST & C & L1 & Blur & 36.55 $\pm$ 1.49 & 0.976 $\pm$ 0.006 \\
MNIST & C & L1 & Noise & 29.48 $\pm$ 0.74 & 0.601 $\pm$ 0.027 \\
MNIST & C & L1 & Resolution & 27.18 $\pm$ 1.51 & 0.927 $\pm$ 0.017 \\
MNIST & C & L1 & S\&P & 19.86 $\pm$ 1.20 & 0.888 $\pm$ 0.022 \\
MNIST & C & L1 & Saturation & 32.35 $\pm$ 1.64 & 0.963 $\pm$ 0.010 \\
MNIST & C & L2 & Blur & 29.67 $\pm$ 1.21 & 0.939 $\pm$ 0.012 \\
MNIST & C & L2 & Noise & 25.36 $\pm$ 0.30 & 0.480 $\pm$ 0.033 \\
MNIST & C & L2 & Resolution & 21.05 $\pm$ 1.58 & 0.847 $\pm$ 0.036 \\
MNIST & C & L2 & S\&P & 16.93 $\pm$ 0.96 & 0.820 $\pm$ 0.030 \\
MNIST & C & L2 & Saturation & 32.35 $\pm$ 1.64 & 0.963 $\pm$ 0.010 \\
MNIST & C & L3 & Blur & 22.34 $\pm$ 1.20 & 0.819 $\pm$ 0.037 \\
MNIST & C & L3 & Noise & 22.23 $\pm$ 0.21 & 0.413 $\pm$ 0.033 \\
MNIST & C & L3 & Resolution & 16.37 $\pm$ 1.23 & 0.730 $\pm$ 0.055 \\
MNIST & C & L3 & S\&P & 14.69 $\pm$ 0.71 & 0.736 $\pm$ 0.033 \\
MNIST & C & L3 & Saturation & 32.36 $\pm$ 1.64 & 0.963 $\pm$ 0.010 \\
MNIST & C & L4 & Blur & 18.18 $\pm$ 1.19 & 0.674 $\pm$ 0.062 \\
MNIST & C & L4 & Noise & 19.86 $\pm$ 0.20 & 0.366 $\pm$ 0.032 \\
MNIST & C & L4 & Resolution & 13.89 $\pm$ 1.19 & 0.645 $\pm$ 0.070 \\
MNIST & C & L4 & S\&P & 13.24 $\pm$ 0.60 & 0.663 $\pm$ 0.034 \\
MNIST & C & L4 & Saturation & 32.36 $\pm$ 1.64 & 0.963 $\pm$ 0.010 \\
MNIST & C & L5 & Blur & 15.50 $\pm$ 1.25 & 0.496 $\pm$ 0.085 \\
MNIST & C & L5 & Noise & 17.16 $\pm$ 0.19 & 0.315 $\pm$ 0.029 \\
MNIST & C & L5 & Resolution & 11.12 $\pm$ 1.80 & 0.507 $\pm$ 0.107 \\
MNIST & C & L5 & S\&P & 12.13 $\pm$ 0.53 & 0.596 $\pm$ 0.035 \\
MNIST & C & L5 & Saturation & 32.36 $\pm$ 1.64 & 0.963 $\pm$ 0.010 \\
MNIST & C2 & L1 & Blur & 36.56 $\pm$ 1.49 & 0.976 $\pm$ 0.006 \\
MNIST & C2 & L1 & Resolution & 27.18 $\pm$ 1.51 & 0.927 $\pm$ 0.017 \\
MNIST & C2 & L1 & S\&P & 19.86 $\pm$ 1.20 & 0.888 $\pm$ 0.022 \\
MNIST & C2 & L2 & Blur & 29.67 $\pm$ 1.21 & 0.939 $\pm$ 0.012 \\
MNIST & C2 & L2 & Resolution & 21.05 $\pm$ 1.58 & 0.847 $\pm$ 0.036 \\
MNIST & C2 & L2 & S\&P & 16.93 $\pm$ 0.96 & 0.820 $\pm$ 0.030 \\
MNIST & C2 & L3 & Blur & 22.34 $\pm$ 1.20 & 0.819 $\pm$ 0.037 \\
MNIST & C2 & L3 & Resolution & 16.37 $\pm$ 1.23 & 0.730 $\pm$ 0.055 \\
MNIST & C2 & L3 & S\&P & 14.69 $\pm$ 0.71 & 0.736 $\pm$ 0.033 \\
MNIST & C2 & L4 & Blur & 18.18 $\pm$ 1.19 & 0.674 $\pm$ 0.062 \\
MNIST & C2 & L4 & Resolution & 13.89 $\pm$ 1.19 & 0.645 $\pm$ 0.070 \\
MNIST & C2 & L4 & S\&P & 13.24 $\pm$ 0.60 & 0.663 $\pm$ 0.034 \\
MNIST & C2 & L5 & Blur & 15.50 $\pm$ 1.25 & 0.496 $\pm$ 0.085 \\
MNIST & C2 & L5 & Resolution & 11.12 $\pm$ 1.80 & 0.507 $\pm$ 0.107 \\
MNIST & C2 & L5 & S\&P & 12.13 $\pm$ 0.53 & 0.596 $\pm$ 0.035 \\
\bottomrule
\end{tabular}
\end{table}


\begin{figure}[htbp]
  \centering
  \includegraphics[width=\textwidth]{\figpath{samples/grid_phase_B_levels_cifar10.png}}
  \caption{CIFAR-10: six fixed validation examples (deterministic seed) at clean, L1 Mild, L3 Moderate, and L5 Extreme combined degradation. Class labels above each column. Rows from top to bottom: Clean, L1, L3, L5.}
  \label{fig:samples_B_cifar}
\end{figure}

\begin{figure}[htbp]
  \centering
  \includegraphics[width=\textwidth]{\figpath{samples/grid_phase_C_axes_L3_cifar10.png}}
  \caption{CIFAR-10: per-axis isolation at L3 (Phase C). Each row shows the same six images with only the named axis active.}
  \label{fig:samples_C_cifar}
\end{figure}

\begin{figure}[htbp]
  \centering
  \includegraphics[width=\textwidth]{\figpath{samples/grid_phase_B_levels_mnist.png}}
  \caption{MNIST: combined degradation at L1, L3, L5 on six fixed validation examples.}
  \label{fig:samples_B_mnist}
\end{figure}

\begin{figure}[htbp]
  \centering
  \includegraphics[width=\textwidth]{\figpath{samples/grid_phase_C_axes_L3_mnist.png}}
  \caption{MNIST: per-axis isolation at L3 (Phase C).}
  \label{fig:samples_C_mnist}
\end{figure}

\subsection{Clean baseline accuracy (Phase A)}
\label{sec:tests:phase_a}

Table \ref{tab:phase_a} reports the validation accuracy at the clean
operating point. These values are the upper bounds against which every
subsequent phase is measured.

\begin{table}[htbp]
\centering
\caption{Phase A clean-baseline accuracy (\%) on the held-out validation split. No degradation applied; values are upper bounds for every subsequent phase.}
\label{tab:phase_a}
\setlength{\tabcolsep}{4pt}
\renewcommand{\arraystretch}{1.1}
\begin{tabular}{lrr}
\toprule
Model & CIFAR-10 & MNIST \\
\midrule
ResNet50 & 95.18 & 99.14 \\
DenseNet121 & 93.56 & 99.30 \\
TransNeXt-tiny & 97.64 & 99.24 \\
\bottomrule
\end{tabular}
\end{table}


\subsection{Combined-degradation accuracy (Phase B)}
\label{sec:tests:phase_b}

Phase B applies every axis simultaneously at the same level. Table
\ref{tab:phase_b} reports the full $3 \times 2 \times 5$ accuracy block;
Figure \ref{fig:acc_vs_level} plots the curves with the Phase D-T3
recovery overlay.

\begin{table}[htbp]
\centering
\caption{Phase B (combined degradation) validation accuracy (\%). All five axes active at the same level L. Cells with $\pm$ values are multi-seed audited (seeds $\{42, 43, 44\}$).}
\label{tab:phase_b}
\setlength{\tabcolsep}{4pt}
\renewcommand{\arraystretch}{1.1}
\begin{tabular}{llrrrrr}
\toprule
Model & Dataset & L1 & L2 & L3 & L4 & L5 \\
\midrule
ResNet50 & CIFAR-10 & 76.46 & 53.94 & 36.31 $\pm$ 0.83 & 26.84 & 19.06 \\
ResNet50 & MNIST & 99.18 & 96.36 & 79.17 $\pm$ 0.80 & 56.44 & 28.00 \\
DenseNet121 & CIFAR-10 & 78.48 & 53.68 & 38.54 $\pm$ 0.78 & 28.52 & 19.58 \\
DenseNet121 & MNIST & 98.74 & 96.54 & 80.79 $\pm$ 0.84 & 58.08 & 27.76 \\
TransNeXt-tiny & CIFAR-10 & 87.58 & 64.22 & 39.48 $\pm$ 0.62 & 28.50 & 20.10 \\
TransNeXt-tiny & MNIST & 98.96 & 96.00 & 80.05 $\pm$ 0.11 & 57.24 & 27.48 \\
\bottomrule
\end{tabular}
\end{table}


\begin{figure}[htbp]
  \centering
  \includegraphics[width=0.49\textwidth]{\figpath{curves/accuracy_vs_level_cifar10.png}}\hfill
  \includegraphics[width=0.49\textwidth]{\figpath{curves/accuracy_vs_level_mnist.png}}
  \caption{Validation accuracy vs combined degradation level. Solid curves: Phase B (no extra regularization). Dashed curves: Phase D-T3 (T3 regularization on the same pipeline). Dotted horizontal lines: Phase A clean baseline per model.}
  \label{fig:acc_vs_level}
\end{figure}

\subsection{Accuracy versus image quality}
\label{sec:tests:psnr}

Figures \ref{fig:acc_vs_psnr} and \ref{fig:acc_vs_ssim} plot validation
accuracy against the per-pipeline mean PSNR and SSIM respectively. The
Spearman rank correlation is reported in each panel and is monotonically
strong on both metrics for both datasets. Accuracy and image quality
collapse together, but the slope is not unity: on CIFAR-10 the model
loses accuracy faster than SSIM, while on MNIST the two metrics track
each other closely. This corroborates the per-axis attribution in
Section \ref{sec:tests:phase_c}: CIFAR-10 carries more inter-class color
information that is destroyed by the saturation lerp, and SSIM does not
penalize the color collapse strongly.

\begin{figure}[htbp]
  \centering
  \includegraphics[width=\textwidth]{\figpath{curves/accuracy_vs_psnr.png}}
  \caption{Validation accuracy vs mean PSNR per pipeline. Each marker is one trained cell; color indicates the backbone. The strong negative monotone relationship justifies treating PSNR as a low-cost proxy for downstream accuracy at design time.}
  \label{fig:acc_vs_psnr}
\end{figure}

\begin{figure}[htbp]
  \centering
  \includegraphics[width=\textwidth]{\figpath{curves/accuracy_vs_ssim.png}}
  \caption{Validation accuracy vs mean SSIM per pipeline.}
  \label{fig:acc_vs_ssim}
\end{figure}

\subsection{Per-axis isolation (Phase C)}
\label{sec:tests:phase_c}

Phase C isolates each axis: the named axis is at level L, every other axis
is at its identity value. Table \ref{tab:phase_c} reports the full
$3 \times 2 \times 5 \times 5 = 150$ accuracy block; Figure
\ref{fig:attribution} shows the per-axis attribution at L3 as
accuracy delta vs the Phase A clean baseline.

\begin{table}[htbp]
\centering
\caption{Phase C single-axis isolation validation accuracy (\%). The named axis is at level L; every other axis is at its identity value.}
\label{tab:phase_c}
\setlength{\tabcolsep}{4pt}
\renewcommand{\arraystretch}{1.1}
\begin{tabular}{lllrrrrr}
\toprule
Model & Dataset & Axis & L1 & L2 & L3 & L4 & L5 \\
\midrule
ResNet50 & CIFAR-10 & Resolution & 88.00 & 78.98 & 68.46 & 60.50 & 43.16 \\
ResNet50 & CIFAR-10 & Blur & 93.54 & 89.88 & 86.78 & 81.46 & 70.98 \\
ResNet50 & CIFAR-10 & Noise & 89.14 & 84.70 & 78.06 & 72.76 & 66.88 \\
ResNet50 & CIFAR-10 & S\&P & 90.42 & 89.92 & 86.52 & 84.64 & 83.30 \\
ResNet50 & CIFAR-10 & Saturation & 92.06 & 92.82 & 92.56 & 91.78 & 88.22 \\
ResNet50 & MNIST & Resolution & 98.84 & 98.50 & 93.48 & 81.22 & 44.26 \\
ResNet50 & MNIST & Blur & 99.32 & 99.14 & 99.18 & 99.20 & 98.82 \\
ResNet50 & MNIST & Noise & 99.10 & 99.16 & 99.10 & 99.10 & 98.96 \\
ResNet50 & MNIST & S\&P & 99.08 & 98.98 & 98.88 & 98.90 & 98.62 \\
ResNet50 & MNIST & Saturation & 98.94 & 99.18 & 99.36 & 99.18 & 99.10 \\
DenseNet121 & CIFAR-10 & Resolution & 86.52 & 79.60 & 67.68 & 57.76 & 44.32 \\
DenseNet121 & CIFAR-10 & Blur & 90.54 & 90.40 & 86.44 & 82.04 & 71.42 \\
DenseNet121 & CIFAR-10 & Noise & 88.92 & 83.64 & 79.54 & 75.06 & 68.46 \\
DenseNet121 & CIFAR-10 & S\&P & 90.72 & 88.98 & 87.96 & 86.44 & 85.08 \\
DenseNet121 & CIFAR-10 & Saturation & 91.84 & 91.42 & 91.66 & 91.30 & 88.34 \\
DenseNet121 & MNIST & Resolution & 99.22 & 98.62 & 93.60 & 81.32 & 45.40 \\
DenseNet121 & MNIST & Blur & 99.06 & 99.34 & 98.96 & 99.28 & 98.46 \\
DenseNet121 & MNIST & Noise & 99.34 & 99.26 & 99.24 & 99.34 & 99.28 \\
DenseNet121 & MNIST & S\&P & 99.18 & 99.14 & 99.06 & 98.76 & 98.72 \\
DenseNet121 & MNIST & Saturation & 99.32 & 99.24 & 99.26 & 99.28 & 99.14 \\
TransNeXt-tiny & CIFAR-10 & Resolution & 95.14 & 88.62 & 79.64 & 68.40 & 46.24 \\
TransNeXt-tiny & CIFAR-10 & Blur & 97.66 & 97.24 & 94.96 & 90.88 & 82.44 \\
TransNeXt-tiny & CIFAR-10 & Noise & 95.98 & 93.34 & 89.40 & 85.98 & 80.08 \\
TransNeXt-tiny & CIFAR-10 & S\&P & 97.28 & 96.60 & 95.48 & 94.74 & 94.22 \\
TransNeXt-tiny & CIFAR-10 & Saturation & 97.92 & 97.78 & 97.54 & 97.56 & 95.34 \\
TransNeXt-tiny & MNIST & Resolution & 99.14 & 98.08 & 92.48 & 80.08 & 44.86 \\
TransNeXt-tiny & MNIST & Blur & 99.24 & 98.96 & 99.14 & 99.28 & 99.06 \\
TransNeXt-tiny & MNIST & Noise & 99.16 & 99.28 & 99.24 & 99.16 & 99.14 \\
TransNeXt-tiny & MNIST & S\&P & 99.00 & 99.16 & 98.82 & 98.76 & 98.94 \\
TransNeXt-tiny & MNIST & Saturation & 99.36 & 98.96 & 99.18 & 99.24 & 99.04 \\
\bottomrule
\end{tabular}
\end{table}


\begin{figure}[htbp]
  \centering
  \includegraphics[width=0.49\textwidth]{\figpath{attribution/axis_attribution_L3_cifar10.png}}\hfill
  \includegraphics[width=0.49\textwidth]{\figpath{attribution/axis_attribution_L3_mnist.png}}
  \caption{Per-axis attribution at L3: accuracy delta vs the Phase A clean baseline. Resolution dominates on CIFAR-10; on MNIST the spread between axes is smaller because the digit shapes are tolerant of color loss and moderate blur.}
  \label{fig:attribution}
\end{figure}

Table \ref{tab:attribution} reports the L3 attribution numerically together
with the combined Phase B L3 cell so the reader can verify that the
per-axis drops sum to less than the combined drop (the axes are not
additive: their interactions compound).

\begin{table}[htbp]
\centering
\caption{Per-axis accuracy attribution at L3: validation-accuracy drop (percentage points) relative to the Phase A clean baseline. "Combined" is the Phase B L3 cell (all five axes active simultaneously).}
\label{tab:attribution}
\setlength{\tabcolsep}{4pt}
\renewcommand{\arraystretch}{1.1}
\begin{tabular}{llrrrrrr}
\toprule
Model & Dataset & Resolution & Blur & Noise & S\&P & Saturation & Combined (B L3) \\
\midrule
ResNet50 & CIFAR-10 & -26.72 & -8.40 & -17.12 & -8.66 & -2.62 & -59.32 \\
ResNet50 & MNIST & -5.66 & +0.04 & -0.04 & -0.26 & +0.22 & -20.28 \\
DenseNet121 & CIFAR-10 & -25.88 & -7.12 & -14.02 & -5.60 & -1.90 & -53.96 \\
DenseNet121 & MNIST & -5.70 & -0.34 & -0.06 & -0.24 & -0.04 & -17.34 \\
TransNeXt-tiny & CIFAR-10 & -18.00 & -2.68 & -8.24 & -2.16 & -0.10 & -58.64 \\
TransNeXt-tiny & MNIST & -6.76 & -0.10 & +0.00 & -0.42 & -0.06 & -19.30 \\
\bottomrule
\end{tabular}
\end{table}


\subsection{Regularization recovery (Phase D)}
\label{sec:tests:phase_d}

Phase D layers three regularization treatments on top of the Phase B Optuna
winners. Table \ref{tab:phase_d} reports the accuracy at every level for
every treatment.

\begin{table}[htbp]
\centering
\caption{Phase D regularization-recovery validation accuracy (\%). T1 = architectural regularization only; T2 = label-mixing only; T3 = combination. Degradation pipeline matches Phase B at the same level.}
\label{tab:phase_d}
\setlength{\tabcolsep}{4pt}
\renewcommand{\arraystretch}{1.1}
\begin{tabular}{lllrrrrr}
\toprule
Model & Dataset & Treatment & L1 & L2 & L3 & L4 & L5 \\
\midrule
ResNet50 & CIFAR-10 & T1 & 76.44 & 52.92 & 35.18 & 27.10 & 18.64 \\
ResNet50 & CIFAR-10 & T2 & 78.12 & 52.76 & 37.50 & 28.22 & 18.34 \\
ResNet50 & CIFAR-10 & T3 & 79.12 & 54.38 & 37.27 $\pm$ 0.43 & 27.88 & 19.24 \\
ResNet50 & MNIST & T1 & 98.50 & 96.16 & 79.70 & 55.62 & 28.14 \\
ResNet50 & MNIST & T2 & 99.12 & 96.24 & 78.58 & 57.00 & 27.34 \\
ResNet50 & MNIST & T3 & 98.98 & 96.96 & 79.85 $\pm$ 0.22 & 56.00 & 27.86 \\
DenseNet121 & CIFAR-10 & T1 & 76.74 & 54.32 & 39.32 & 29.64 & 20.12 \\
DenseNet121 & CIFAR-10 & T2 & 76.64 & 55.76 & 39.16 & 29.82 & 19.72 \\
DenseNet121 & CIFAR-10 & T3 & 77.04 & 54.80 & 39.56 $\pm$ 0.23 & 29.96 & 19.90 \\
DenseNet121 & MNIST & T1 & 98.90 & 96.82 & 80.74 & 57.56 & 27.96 \\
DenseNet121 & MNIST & T2 & 99.06 & 96.18 & 79.84 & 57.78 & 28.02 \\
DenseNet121 & MNIST & T3 & 99.30 & 97.02 & 81.49 $\pm$ 0.22 & 57.58 & 27.26 \\
TransNeXt-tiny & CIFAR-10 & T1 & 88.46 & 66.96 & 42.36 & 28.44 & 19.66 \\
TransNeXt-tiny & CIFAR-10 & T2 & 89.42 & 65.26 & 42.34 & 29.88 & 20.02 \\
TransNeXt-tiny & CIFAR-10 & T3 & 90.28 & 68.00 & 43.33 $\pm$ 0.96 & 29.72 & 20.42 \\
TransNeXt-tiny & MNIST & T1 & 98.92 & 96.04 & 80.22 & 58.94 & 27.62 \\
TransNeXt-tiny & MNIST & T2 & 99.18 & 97.02 & 80.42 & 55.96 & 27.52 \\
TransNeXt-tiny & MNIST & T3 & 99.38 & 97.10 & 81.20 $\pm$ 0.50 & 59.24 & 27.94 \\
\bottomrule
\end{tabular}
\end{table}


\begin{figure}[htbp]
  \centering
  \includegraphics[width=0.85\textwidth]{\figpath{phase_d/recovery_summary.png}}
  \caption{Phase D recovery summary. Each bar reports the mean accuracy delta vs the Phase B v2 baseline at the same level, averaged across the five levels. T3 is the best treatment but the absolute recovery remains modest (mean $\sim 0.77$ \Mpp{}) compared with the Phase B v2 collapse vs v1 ($-12.49$ \Mpp{}).}
  \label{fig:phase_d_recovery}
\end{figure}

\subsection{THz protocol simplification (Phase B2 and B2-nr)}
\label{sec:tests:phase_b2}

Phase B2 simplifies the THz protocol by forcing saturation $=0$ (full
grayscale) and noise $=0$ (no additive noise). T3 regularization is still
applied. Phase B2-nr is the no-regularization arm at L3 only, used to
decompose the Phase B1 to B2 improvement into pure-protocol and
regularization fractions. Table \ref{tab:phase_b2} reports the accuracy
block.

\begin{table}[htbp]
\centering
\caption{Phase B2 (THz-protocol simplification, saturation $=0$, noise $=0$) and Phase B2-nr (L3 only, no regularization) validation accuracy (\%).}
\label{tab:phase_b2}
\setlength{\tabcolsep}{4pt}
\renewcommand{\arraystretch}{1.1}
\begin{tabular}{lllrrrrr}
\toprule
Model & Dataset & Arm & L1 & L2 & L3 & L4 & L5 \\
\midrule
ResNet50 & CIFAR-10 & B2 (T3) & 74.96 & 55.32 & 41.09 $\pm$ 0.45 & 32.66 & 24.86 \\
ResNet50 & CIFAR-10 & B2nr & -- & -- & 40.31 $\pm$ 0.41 & -- & -- \\
ResNet50 & MNIST & B2 (T3) & 98.90 & 97.14 & 84.88 $\pm$ 0.49 & 63.00 & 34.84 \\
ResNet50 & MNIST & B2nr & -- & -- & 85.74 $\pm$ 0.04 & -- & -- \\
DenseNet121 & CIFAR-10 & B2 (T3) & 73.98 & 56.68 & 42.41 $\pm$ 0.93 & 34.70 & 24.88 \\
DenseNet121 & CIFAR-10 & B2nr & -- & -- & 40.75 $\pm$ 0.32 & -- & -- \\
DenseNet121 & MNIST & B2 (T3) & 98.86 & 97.62 & 86.23 $\pm$ 0.18 & 66.30 & 35.22 \\
DenseNet121 & MNIST & B2nr & -- & -- & 85.70 $\pm$ 0.56 & -- & -- \\
TransNeXt-tiny & CIFAR-10 & B2 (T3) & 86.52 & 68.20 & 46.43 $\pm$ 0.02 & 36.46 & 26.50 \\
TransNeXt-tiny & CIFAR-10 & B2nr & -- & -- & 43.57 $\pm$ 0.27 & -- & -- \\
TransNeXt-tiny & MNIST & B2 (T3) & 99.18 & 97.62 & 85.67 $\pm$ 0.25 & 65.24 & 34.94 \\
TransNeXt-tiny & MNIST & B2nr & -- & -- & 84.61 $\pm$ 0.28 & -- & -- \\
\bottomrule
\end{tabular}
\end{table}


\begin{figure}[htbp]
  \centering
  \includegraphics[width=0.85\textwidth]{\figpath{phase_b2/comparison_summary.png}}
  \caption{Phase B2 protocol-simplification decomposition summary. The B1 to B2 accuracy gain ($\sim +4.01$ \Mpp{} on average) decomposes as $\sim +3.24$ \Mpp{} pure-protocol effect (saturation and noise off) and $\sim +0.77$ \Mpp{} regularization effect.}
  \label{fig:phase_b2_summary}
\end{figure}

\subsection{Single-axis under the THz protocol (Phase C2)}
\label{sec:tests:phase_c2}

Phase C2 isolates each of the three THz-relevant axes (resolution, blur,
salt and pepper) under the simplified protocol. Table \ref{tab:phase_c2}
reports the accuracy block; Figure \ref{fig:phase_c2_summary} shows the
attribution summary.

\begin{table}[htbp]
\centering
\caption{Phase C2 single-axis isolation under the THz protocol (saturation $=0$, noise $=0$, T3 regularization) validation accuracy (\%).}
\label{tab:phase_c2}
\setlength{\tabcolsep}{4pt}
\renewcommand{\arraystretch}{1.1}
\begin{tabular}{lllrrrrr}
\toprule
Model & Dataset & Axis & L1 & L2 & L3 & L4 & L5 \\
\midrule
ResNet50 & CIFAR-10 & Resolution & 80.18 & 66.36 & 51.32 & 42.96 & 30.16 \\
ResNet50 & CIFAR-10 & Blur & 90.62 & 89.58 & 83.36 & 71.68 & 58.76 \\
ResNet50 & CIFAR-10 & S\&P & 86.24 & 85.18 & 83.48 & 77.86 & 74.68 \\
ResNet50 & MNIST & Resolution & 99.06 & 98.64 & 94.12 & 81.24 & 45.54 \\
ResNet50 & MNIST & Blur & 99.24 & 99.26 & 99.32 & 99.32 & 98.86 \\
ResNet50 & MNIST & S\&P & 99.14 & 99.16 & 99.18 & 99.00 & 99.00 \\
DenseNet121 & CIFAR-10 & Resolution & 76.82 & 64.04 & 51.34 & 44.44 & 30.80 \\
DenseNet121 & CIFAR-10 & Blur & 88.24 & 85.66 & 81.66 & 74.90 & 64.60 \\
DenseNet121 & CIFAR-10 & S\&P & 86.38 & 83.26 & 80.54 & 79.30 & 77.78 \\
DenseNet121 & MNIST & Resolution & 99.38 & 98.64 & 94.30 & 82.28 & 45.02 \\
DenseNet121 & MNIST & Blur & 99.42 & 99.34 & 99.22 & 99.38 & 99.16 \\
DenseNet121 & MNIST & S\&P & 99.16 & 99.20 & 99.20 & 99.08 & 99.06 \\
TransNeXt-tiny & CIFAR-10 & Resolution & 89.34 & 76.08 & 58.24 & 46.58 & 30.54 \\
TransNeXt-tiny & CIFAR-10 & Blur & 96.52 & 95.96 & 93.00 & 88.02 & 71.04 \\
TransNeXt-tiny & CIFAR-10 & S\&P & 95.20 & 94.78 & 93.52 & 91.56 & 89.58 \\
TransNeXt-tiny & MNIST & Resolution & 99.32 & 98.76 & 93.94 & 81.74 & 45.60 \\
TransNeXt-tiny & MNIST & Blur & 99.38 & 99.38 & 99.46 & 99.48 & 99.06 \\
TransNeXt-tiny & MNIST & S\&P & 99.28 & 99.14 & 99.20 & 99.08 & 99.08 \\
\bottomrule
\end{tabular}
\end{table}


\begin{figure}[htbp]
  \centering
  \includegraphics[width=0.85\textwidth]{\figpath{phase_c2/attribution_summary.png}}
  \caption{Phase C2 per-axis THz attribution summary. Mean accuracy delta from L1 to L5: resolution $-27.45$ \Mpp{}, blur $-6.58$ \Mpp{}, salt and pepper $-5.13$ \Mpp{}. Resolution is the dominant axis by a factor of approximately four.}
  \label{fig:phase_c2_summary}
\end{figure}

\subsection{Multi-seed audit}
\label{sec:tests:multiseed}

The multi-seed audit re-runs each of the 24 L3 headline cells at
seeds 43 and 44 in addition to the canonical seed 42. Table
\ref{tab:multiseed} reports the resulting mean $\pm \sigma$ values. Every
cell has $\sigma < 1.5$ \Mpp{}, which bounds the residual run-to-run
variance in every comparison reported above.

\begin{table}[htbp]
\centering
\caption{Multi-seed audit: validation accuracy at L3 across seeds $\{42, 43, 44\}$, on the 24 L3 headline cells. Standard deviation is reported in percentage points.}
\label{tab:multiseed}
\setlength{\tabcolsep}{4pt}
\renewcommand{\arraystretch}{1.1}
\begin{tabular}{lllllrrr}
\toprule
Phase & Model & Dataset & Level & Axis / Tr. & Mean (\%) & $\sigma$ (pp) & $n$ \\
\midrule
B & DenseNet121 & CIFAR-10 & L3 & -- & 38.54 & 0.78 & 3 \\
B & DenseNet121 & MNIST & L3 & -- & 80.79 & 0.84 & 3 \\
B & ResNet50 & CIFAR-10 & L3 & -- & 36.31 & 0.83 & 3 \\
B & ResNet50 & MNIST & L3 & -- & 79.17 & 0.80 & 3 \\
B & TransNeXt-tiny & CIFAR-10 & L3 & -- & 39.48 & 0.62 & 3 \\
B & TransNeXt-tiny & MNIST & L3 & -- & 80.05 & 0.11 & 3 \\
B2 & DenseNet121 & CIFAR-10 & L3 & T3 & 42.41 & 0.93 & 3 \\
B2 & DenseNet121 & MNIST & L3 & T3 & 86.23 & 0.18 & 3 \\
B2 & ResNet50 & CIFAR-10 & L3 & T3 & 41.09 & 0.45 & 3 \\
B2 & ResNet50 & MNIST & L3 & T3 & 84.88 & 0.49 & 3 \\
B2 & TransNeXt-tiny & CIFAR-10 & L3 & T3 & 46.43 & 0.02 & 3 \\
B2 & TransNeXt-tiny & MNIST & L3 & T3 & 85.67 & 0.25 & 3 \\
B2nr & DenseNet121 & CIFAR-10 & L3 & -- & 40.75 & 0.32 & 3 \\
B2nr & DenseNet121 & MNIST & L3 & -- & 85.70 & 0.56 & 3 \\
B2nr & ResNet50 & CIFAR-10 & L3 & -- & 40.31 & 0.41 & 3 \\
B2nr & ResNet50 & MNIST & L3 & -- & 85.74 & 0.04 & 3 \\
B2nr & TransNeXt-tiny & CIFAR-10 & L3 & -- & 43.57 & 0.27 & 3 \\
B2nr & TransNeXt-tiny & MNIST & L3 & -- & 84.61 & 0.28 & 3 \\
D & DenseNet121 & CIFAR-10 & L3 & T3 & 39.56 & 0.23 & 3 \\
D & DenseNet121 & MNIST & L3 & T3 & 81.49 & 0.22 & 3 \\
D & ResNet50 & CIFAR-10 & L3 & T3 & 37.27 & 0.43 & 3 \\
D & ResNet50 & MNIST & L3 & T3 & 79.85 & 0.22 & 3 \\
D & TransNeXt-tiny & CIFAR-10 & L3 & T3 & 43.33 & 0.96 & 3 \\
D & TransNeXt-tiny & MNIST & L3 & T3 & 81.20 & 0.50 & 3 \\
\bottomrule
\end{tabular}
\end{table}


\begin{figure}[htbp]
  \centering
  \includegraphics[width=0.49\textwidth]{\figpath{multiseed/L3_variance_cifar10.png}}\hfill
  \includegraphics[width=0.49\textwidth]{\figpath{multiseed/L3_variance_mnist.png}}
  \caption{Multi-seed variance at L3: each bar reports the mean validation accuracy across seeds $\{42, 43, 44\}$ with the standard-deviation error bar. Color encodes the phase. The $\sigma$ bars are small relative to the cross-phase deltas, which validates the cross-phase comparisons.}
  \label{fig:multiseed_variance}
\end{figure}

\subsection{Calibration}
\label{sec:tests:calibration}

Figure \ref{fig:reliability_L3} shows the reliability diagrams of the six
Phase B L3 headline cells; Figure \ref{fig:reliability_L5} shows the same
cells at L5. The diagonal would correspond to perfect calibration. Both
CNN backbones inflate confidence as the spatial information collapses,
which is consistent with the reported ECE growth from below 5 \Mpp{} at
L3 to above 15 \Mpp{} at L5 on CIFAR-10 \cite{guo2017oncalibration}.

\begin{figure}[htbp]
  \centering
  \begin{minipage}{0.32\textwidth}\includegraphics[width=\linewidth]{\figpath{calibration/B_L3_resnet50_cifar10.png}}\end{minipage}\hfill
  \begin{minipage}{0.32\textwidth}\includegraphics[width=\linewidth]{\figpath{calibration/B_L3_densenet121_cifar10.png}}\end{minipage}\hfill
  \begin{minipage}{0.32\textwidth}\includegraphics[width=\linewidth]{\figpath{calibration/B_L3_transnext_tiny_cifar10.png}}\end{minipage}\\[0.3em]
  \begin{minipage}{0.32\textwidth}\includegraphics[width=\linewidth]{\figpath{calibration/B_L3_resnet50_mnist.png}}\end{minipage}\hfill
  \begin{minipage}{0.32\textwidth}\includegraphics[width=\linewidth]{\figpath{calibration/B_L3_densenet121_mnist.png}}\end{minipage}\hfill
  \begin{minipage}{0.32\textwidth}\includegraphics[width=\linewidth]{\figpath{calibration/B_L3_transnext_tiny_mnist.png}}\end{minipage}
  \caption{Reliability diagrams at Phase B L3. Top row: CIFAR-10. Bottom row: MNIST. Columns: ResNet50, DenseNet121, TransNeXt-tiny.}
  \label{fig:reliability_L3}
\end{figure}

\begin{figure}[htbp]
  \centering
  \begin{minipage}{0.32\textwidth}\includegraphics[width=\linewidth]{\figpath{calibration/B_L5_resnet50_cifar10.png}}\end{minipage}\hfill
  \begin{minipage}{0.32\textwidth}\includegraphics[width=\linewidth]{\figpath{calibration/B_L5_densenet121_cifar10.png}}\end{minipage}\hfill
  \begin{minipage}{0.32\textwidth}\includegraphics[width=\linewidth]{\figpath{calibration/B_L5_transnext_tiny_cifar10.png}}\end{minipage}\\[0.3em]
  \begin{minipage}{0.32\textwidth}\includegraphics[width=\linewidth]{\figpath{calibration/B_L5_resnet50_mnist.png}}\end{minipage}\hfill
  \begin{minipage}{0.32\textwidth}\includegraphics[width=\linewidth]{\figpath{calibration/B_L5_densenet121_mnist.png}}\end{minipage}\hfill
  \begin{minipage}{0.32\textwidth}\includegraphics[width=\linewidth]{\figpath{calibration/B_L5_transnext_tiny_mnist.png}}\end{minipage}
  \caption{Reliability diagrams at Phase B L5. Top row: CIFAR-10. Bottom row: MNIST. Columns: ResNet50, DenseNet121, TransNeXt-tiny.}
  \label{fig:reliability_L5}
\end{figure}

\subsection{Confusion matrices at L5}
\label{sec:tests:confusion}

Figure \ref{fig:confusion_L5} shows the six confusion matrices at the
Phase B L5 cell, the worst case across the campaign. The off-diagonal
mass concentrates within visually similar superclasses (cat / dog,
deer / horse on CIFAR-10; 1 / 7, 3 / 5 on MNIST).

\begin{figure}[htbp]
  \centering
  \begin{minipage}{0.32\textwidth}\includegraphics[width=\linewidth]{\figpath{confusion/B_L5_resnet50_cifar10.png}}\end{minipage}\hfill
  \begin{minipage}{0.32\textwidth}\includegraphics[width=\linewidth]{\figpath{confusion/B_L5_densenet121_cifar10.png}}\end{minipage}\hfill
  \begin{minipage}{0.32\textwidth}\includegraphics[width=\linewidth]{\figpath{confusion/B_L5_transnext_tiny_cifar10.png}}\end{minipage}\\[0.3em]
  \begin{minipage}{0.32\textwidth}\includegraphics[width=\linewidth]{\figpath{confusion/B_L5_resnet50_mnist.png}}\end{minipage}\hfill
  \begin{minipage}{0.32\textwidth}\includegraphics[width=\linewidth]{\figpath{confusion/B_L5_densenet121_mnist.png}}\end{minipage}\hfill
  \begin{minipage}{0.32\textwidth}\includegraphics[width=\linewidth]{\figpath{confusion/B_L5_transnext_tiny_mnist.png}}\end{minipage}
  \caption{Confusion matrices at Phase B L5. Top row: CIFAR-10. Bottom row: MNIST. Columns: ResNet50, DenseNet121, TransNeXt-tiny.}
  \label{fig:confusion_L5}
\end{figure}

\subsection{Inference throughput}
\label{sec:tests:throughput}

Figure \ref{fig:throughput} reports the wall-clock inference latency for
each backbone on the test rig. ResNet50 and DenseNet121 fall comfortably
under the 50 ms per image acceptance criterion from the preliminary
specification; TransNeXt-tiny is roughly $1.5\times$ slower because
\texttt{torch.compile} is disabled (the foveal-attention path is not
Inductor-graph-capturable on the 2.11 stack).

\begin{figure}[htbp]
  \centering
  \includegraphics[width=0.65\textwidth]{\figpath{inference_throughput.png}}
  \caption{Inference latency per image at $224 \times 224$, batch size 1, \texttt{bf16-mixed}, RTX 5070.}
  \label{fig:throughput}
\end{figure}
